\chapter{Elliptic curves}
Taken from \emph{A Graduate Course in Applied Cryptography, Dan Boneh and Victor Shoup}

The best-known discrete log algorithm in an elliptic curve group of size $q$ runs in time $O(\sqrt{q})$. 
This means that to provide security comparable to AES-128, it suffices to use a group of size $q \approx 2^{256}$ 
so that the time to compute discrete log is $\sqrt{q} \approx 2^{128}$.

\section{a bit of history on elliptic curves}
\sloppy{
Elliptic curves come up naturally in several branches of mathematics.
Here we will follow their development as a branch of arithmetic (the study of rational numbers).
Diophantus, a greek mathematician who lived in Alexandria in the third century, was interested 
in the following problem: Given a bivariate polynomial equation, $f(x,y)=0$ find 
rational points satisfying the equation. That is, find $P=(x,y)$ such that $x,y\in\mathbb{Q}$.
Diophantus wrote a series of influential books on this subject, called the \emph{Arithmetica}, of which six survived.
Much of Arithmathica studies integer and rational solutions of quadratic equations.
For example, one of the first problems in the books regarding elliptic curves is to find $(x,y) \in \Q$
 that satisfy the equation $$ y^{2}=x^3 - x + 9$$
The method invented to answer this question secures most Internet traffic worldwide.
One can easily verify that the following points are on the curve $(0, \pm3), (1, \pm3), (-1, \pm 3)$.
Diophantus developed a method similar to \underline{chord method}:
Given two rational points $U$ and $V$ on the curve, where $U\neq -V$, we can pass a line through them,
 and this line must intersect the curve at a third rational point $W$. 
The chord method was rediscovered several times over the centuries, but it finally stuck with the work of Poincare.

Poincare defines the sum of $U, V$, denoted $U \boxplus V$ as $U \boxplus V := -W$
Diophantus' method, the \emph{tangent method} is another way to build a new rational point from a given rational point. 
We should later see that the method corresponds to adding the point $P$ to itself: $P\boxplus P$.}

\section{elliptic curves over finite fields}
We are primarily interested in elliptic curves over finite fields for cryptographic applications.
\begin{definition}
  Let $p >3 $ be a prime. An \textbf{elliptic curve} $E$ defined over $\mathbb{F}_p$ is an 
  equation  
  $$ y^2 = x^3 +ax + b$$
  where $a,b \in \mathbb{F}_p$ satisfy $4a^3 +27b^2 \neq 0$. We write $E/\mathbb{F}_p$ to denote
  the fact that $E$ is defined over $\F_p$.
\end{definition}
The weird condition $4a^3 +27b^2 \neq 0 $ ensures that the equation 
$ x^3 + ax + b = 0$ does not have a double root (needed to avoid degeneracies).

\begin{definition}
  We say that point $P=(x,y)$, where $x,y \in \F_p$, is on curve $E$ if it satisfies the equation of $E$.
\end{definition}
\begin{definition}
  $E(\F)$ denotes the set of all points on the curve E and are defined over $\F$.
\end{definition}

The addition law, and adding to the set $E(F_p)$ the member $\mathcal{O}$,
 which servers as the identity member, turns it into a group. I'm not getting into the addition laws.

\section{Elliptic curve cryptography}


Given two points $P$ and $\alpha \cdot P$, where $\alpha \in \Z_q$, it is hard to compute $\alpha$.
The best known algorithm to compute $\alpha$ is $\Omega(\sqrt{q})$.
There are a few exceptions to these rules, and thus one should be careful when defining a new curve.
So it's better to use some specific already implemented and used curve rather than 
generating a random prime $p$ and a random curve over $\F_p$.
\todo[inline]{TODO, put here an encryption scheme, and a signature scheme (can use schnorr here honestly)
The most important thing to remember is that elliptic curve cryptography 
assumes hard exponent, and everything is similar to those algorithms
}

\subsection{ECDSA}
The ECDSA is essential DSA over elliptic curves.
The most commonly known is the signature algorithm used by bitcoin and other cryptocurrencies:
ECDSA, which commonly refers to the following algorithm:

\begin{description}
  \item[Key generation] Grab a random number $x\in \mathbb{F}_q$ where the elliptic curve  
  is a vector space over $\mathbb{F}_q$, and $q$ is the order of the group.
  and return $(sk,pk)=(x,xG)$, where G is the generator of the elliptic curve 
  (Usualy $G$ is a specific base point, because over a a group with $q$ points, 
  almost every point is a generator - due to Fermat's little theorm \ref{fermats-little-theorem}).
  \item[Signing] $(r,s)=\sigma = Sig_{sk}(m)$ is generated as follows:
  \begin{enumerate}
    \item generate a random scalar $k\in \mathbb{F}_q$.
    \item compute $R=g^{k^{-1}}$
    \item compute $r=H(R)$ (notice that this is now a scalar).
    If $r mod=0$ then go back to step 1.
    \item compute $s=k^{-1}\cdot(H(m)+x\cdot r)$\todo{This part is incorrect because
     they want specific bits from $H(m)$, but it's not crucial to understand the algorithm,
      so I skipped it.}
    \item set $\sigma=(r,s)$.
  \end{enumerate}
  \item[Verification] To verify a signature $\sigma=(r,s)$ on message $m$ with 
  public key $pk$, one needs to compute the following 
  $ R'= G^{m\cdot s^{-1}}\cdot pk^{r\cdot s^-1} $ (note that this is a point in the EC).
  Notice that if you subtitute $s$ and $pk$ to their actual values, you should get $R$ if 
  the signature is valid.
  Then, check that $r=H(R')$, if it is accept the signature.
\end{description}

\subsection{Schnorr Signature}
\todo[inline]{need to fill this...}
\section{Pairing based cryptography}
We now show that certain elliptic curves have an additional structure,
called a pairing, which enables a world of new schemes that could not be built
 from discrete log groups without this other structure.

 \begin{definition}
  Let $\G_0, \G_1, \G_T$ be three cyclic groups of prime order $q$ where 
  $g_0\in\G_0$ and $g_1\in\G_1$ are generators. A \textbf{pairing} is an efficiently
  computable function $e: \G_0 \times \G_1 \to \G_T$ satisfying the 
  following properties:
  \begin{itemize}
    \item bilinear: $\forall u,u'\in\G_0, v,v'\in\G_1$  we have 
        $$ e(u\cdot u',v) = e(u,v) \cdot e(u',v) and  e(u,v\cdot v') = e(u,v)\cdot e(u,v')$$
    \item non-degenerate: $g_T := e(g_0,g_1)$  is a generator of $\G_T$.
  \end{itemize}
\end{definition}
We refer to $\G_0,\G_1$ as the pairing groups and $\G_T$ as the target group.

Bilinearity implies the following central property of pairing that will be used
in all our constructions: for all $a, b\in \Z_q$ we have
  $$ e(g_0^a,g_1^b)=e(g_0,g_1)^{a\cdot b}= e(g_0^b,g_1^a)$$
This equality follows from the equalities 
$e(g_0^a,g_1^b)=e(g_0,g_1^b)^a=e(g_0^a,g_1^b)=e(g_0^a,g_1)^b$
which are themselves a direct consequence of bilinearity (Note that in cyclic groups
 such as $\G_0 and \G_1$ this definition is equivalent to the definition of bilinearity above).
[if I understand correctly, due to $\G_0 and \G_1$ having generators, we can represent every number
as a power of the generators themselves, and using bilinearity we can reach the above equations]

In a symmetric pairing group, where $G_0 = G_1$ schemes like ElGamal aren't secure anymore,
that is, given two public keys, it is to break the DDH assumption: 
  $$ e(g^x,g^y) = e(g,g)^{x\cdot y}$$ 
Hence, it is possible to extract information. 

In asymmetric pairing, where $G_0 \neq G_1$, it may still be possible
that the DDH assumption holds in both $G_0$ and $G_1$. 

Furthermore, if the discrete log in $G_T$ is easy, it is easy in $G_0$ (and $G_1$).

\begin{rem}
  The pairing function $e$ on an elliptic curve comes from an algebraic pairing called the \textbf{Weil pairing}.
  Which can be efficiently evaluated. In practice, one uses variants of the weil 
  pairing, called the Tate and Ate pairings.
\end{rem}

We'll avoid $bn256$ curve as it isn't secure.


\subsection*{optimisations}
\begin{itemize}
  \item When the first value of the pairing is known, we can speed things up
  by calculating a table of the other possible values. (though it seems very inefficient, and memory intensive.)
\end{itemize}

\section{Advanced encryption schemes from pairings}

\subsection{The BLS signature scheme}
Let $e:F_0 \times G_1 \to G_T$ be a pairing where $G_0,G_1,G_T$ are
cyclic groups of prime order $q$, and where $g_0\in G_0, g_1\in G_1$ are
generators. 
We'll also require a hash function $H : \mathcal{M} \to G_0$.

The BLS signature scheme denoted $\mathcal{S}_BLS = (G,S,V)$, has message space
$\mathcal{M}$ and works as follows:
\begin{itemize}
  \item $G():$ The key generation algorithm runs as follows 
  $$  \alpha \from \Z_q,u\from g_1^\alpha \in G_1  $$
  The public key is $ pk := u$, and the secret key is $\alpha$.

  \item $S(sk,m):$ To sign a message $m\in \mathcal{M}$ using a secret key
  $sk=\alpha\in\Z_q$ do:
  $$ \sigma\from H(m)^{\alpha} \text{ output } \sigma$$

  \item $V(pk,m,\sigma):$ To verify a signature 
  $$ e(H(m), u) = e(\sigma,g_1) $$
  that is, the signature is accepted due to the following calculations
  $$e(H(m), g_1^\alpha) = e(H(m),g_1)^\alpha = e(H(m)^\alpha,g_1)= 
  e(\sigma, g_1)
  $$
\end{itemize}

\subsection{Identity based encryption}
Using Identity-Based Encryption (IBE from now), Alice can send an encrypted message to Bob
without asking for bobs key beforehand. Beyond key exchange, IBE can also 
be used to construct CPA-secure public key encryption and even search encrypted data.
In an IBE scheme, a trusted entity, Trudy, generates a master key pair $(mpk, msk)$.
The key $mpk$ is the \emph{master public key} and is known to everyone. Trudy keeps the \emph{master 
secret key} $msk$ to herself. When Bob wants to use his email address as his public key, he must
somehow obtain the private key derived from $msk$ and Bob's public key.
We denote Bob's email address as his (public) identity $id$ and denote the corresponding private key 
by $sk_{id}$. He receives $sk_{id}$ by contacting Trudy, she'll use $msk$ and $id$ to generate $sk_{id}$,
and send it back to Bob.
Now anyone, including Alice, can send bob messages without looking up his public key. 
Assuming, of course, that Alice knows $mpk$.

\begin{definition}
  An \textbf{Identiy based encryption scheme $ \mathcal{E}_id =(S,G,E,D)$}
  is a tuple of four efficient algorithms: a \textbf{setup algorithm $S$}, a 
  \textbf{Key generation algorithm $G$}, an \textbf{encryption algorithm $E$}
  and a \textbf{decryption algorithm $D$}.

  \begin{itemize}
    \item $S$ is a probabilistic algorithm which outputs $(mpk, msk) \overset{_R}{\leftarrow} S()$.
    \item $G$ is a probabilistic algorithm invoked as $sk_{id} \getrandom G(msk, id)$.
    \item $E$ is a \emph{probabilistic} algorithm invoked as $c \getrandom E(mpk, id, m)$.
    \item $D$ is a deterministic algorithm invoked as $m\getrandom D(sk_{id},c)$. Where $m$ is
      either a message or a \textbf{reject} value.
    \item We require that decryption undoes encryption: 
    $$ \Pr[D(G(msk,id), E(mpk,id,m))=m]=1$$
    \item identities, messages, and ciphertexts are in their respective finite spaces,
    and $\mathcal{E}_id$ is defined over $(\mathcal{ID},\mathcal{M},\mathcal{C})$.
  \end{itemize}
\end{definition}

We can view IBE as a particular public key encryption scheme where the messages
to be encrypted pairs $(id, m)\in\mathcal{ID\times M}$. The master secret key can
decrypt any well-formed ciphertext. However, there are weaker secret keys, 
such as $sk_{id}$ that can decrypt a ciphertext $c\getrandom E(mpk, id', m)$ only if $id'=id$.


Bob won't use his email address as a public key if a company uses IBE. Otherwise, if Bob's secret key 
is compromised, bob would need to change his email. Instead, Bob's public key
should bob could use the identity string $(id, date)$.
Bob will need to request a new secret key from Trudy every day.
Now Alice can even encrypt emails for specific dates, and bob cannot access it before that, 
unless Trudy gives bob any key he requests.

I will not go into detail about the formal definition of security, but the idea is that
the attacker has access to any $sk_{id'}$ it want but should not be able to break semantic
security for some other identity with secret key $sk_{id}$ that the adversary does not have.


\subsection*{IBE from pairings}
We'll need a few components before being able to construct an IBE: 
\begin{itemize}
  \item a pairing. as defined above, over cyclic groups of prime order $q$, 
  with generators $ g_0, g_1 $.
  \item a symmetric cipher $\mathcal{E}=(E_s,D_s)$.
  \item hash functions $H_0:\mathcal{ID}\to \G_0$ and $H_1:\G_1\times\G_T\to \mathcal{k}$,
  where k is the key space.
\end{itemize}
\subsubsection*{Construction 1}

\begin{itemize}
  \item $S():$ runs as follows:
  $$ \alpha\getrandom \Z_q, u_1\from g^\alpha, mpk \from u_1, msk \from \alpha 
  \text{ output }(mpk,msk) $$
  \item $G(msk,id):$ key generation using $msk = \alpha$:
  $$ sk_id \from H_0(id)^\alpha \in \G_0 \text{ output } sk_id$$
  \item $E(mpk, id, m ):$ encryption using the public parameters $mpk=u_1$:
  \begin{align*}\label{lamlam} 
    \beta \getrandom \Z_q,\  z\from e(H_0(id),mpk^\beta) 
    \\ \ 
    k\from H_1(g^\beta,z), c\getrandom E_s(k,m), \text{ output } (g^\beta,c)
  \end{align*}
  The encryption generates a new secret key $k$ at random. This private key relies on the pairing with
  the hashed public id. Once we've generated the ephemeral secret key $k$, we encrypt the message $m$.

  \item $D(sk_id,(w_1,c)):$ decryption using secret key $sk_id$ of ciphertext $(w_1,c)$:
  \begin{align*}
    z \from  e(sk_{id},w_1), k\from H_1(w_1, z), m\from D_s(k,c),
  \text{ output } m
  \end{align*}
  The decryption relies on $w_1=g^\beta$ to figure out the symmetric encryption key $k$.
  With $k$, the message can be decrypted.
  \newline
  To generate $k$ correctly, we first remind the reader that $$sk=H(id)^\alpha, mpk^\beta=g^{\alpha\beta}, k=H(g^\beta,z)$$
  now, let's compute $z$, our missing component:
  $$
  z=\underbrace{e(sk_{id},g^\beta)}_{\text{z in decryption}}=e(H(id)^\alpha,g^\beta)=e(H(id),g^{\alpha\beta})=\underbrace{e(H(id),mpk^\beta)}_{\text{z in encryption}}
  $$
\end{itemize}
